\documentclass[../../../include/open-logic-section]{subfiles}

\begin{document}

\olfileid[fr]{sth}{card-arithmetic}{opps}
\olsection{Définition des opérations fondamentales}

Comme nous n'avons pas à tenir compte de l'ordre,
l'arithmétique cardinale se définit plus facilement que
l'arithmétique ordinale. Nous allons définir simultanément
l'addition, la multiplication et l'exponentiation.

\begin{defn}
Si $\cardfont{a}$ et $\cardfont{b}$ sont des cardinaux :
\begin{align*}
	\cardfont{a} \cardplus \cardfont{b} &\defis
	\card{\cardfont{a} \disjointsum \cardfont{b}}\\
	\cardfont{a} \cardtimes \cardfont{b} &\defis
	\card{\cardfont{a} \times \cardfont{b}}\\
	\cardexpo{\cardfont{a}}{\cardfont{b}} &\defis
	\card{\funfromto{\cardfont{b}}{\cardfont{a}}}
\end{align*}
où $\funfromto{X}{Y} = \Setabs{f}{f \text{ est une fonction } X \to Y}$.
(Il est facile de montrer que $\funfromto{X}{Y}$ existe
pour tous ensembles $X$ et $Y$ ; nous laissons cela en exercice.)
\end{defn}

\begin{prob}
Démontrez dans $\Zminus$ que $\funfromto{X}{Y}$ existe
pour tous ensembles $X$ et~$Y$. Dans $\ZF$, calculez
$\setrank{\funfromto{X}{Y}}$ à partir de
$\setrank{X}$ et $\setrank{Y}$, sur le modèle du
\olref[sth][ord-arithmetic][using-addition]{rankcomputation}.
\end{prob}

Quelques explications peuvent éclairer cette définition.
Pour l'addition, nous utilisons la somme disjointe
$\disjointsum$, définie dans la
\olref[ord-arithmetic][add]{defdissum} ; on voit
facilement que notre définition donne le bon résultat
pour les cas finis. Pour la multiplication,
la \olref[sfr][set][pai]{cardnmprod} nous dit que,
si $A$ possède $n$ éléments et $B$ en possède $m$,
alors $A \times B$ possède $n \cdot m$ éléments :
notre définition étend simplement cette idée à la
multiplication transfinie. Il en va de même pour
l'exponentiation : nous étendons l'idée du cas fini
au cas transfini. Sous certains aspects, l'arithmétique
cardinale transfinie ressemble même beaucoup plus à
l'arithmétique «~ordinaire~» que ne le fait
l'arithmétique ordinale transfinie :

\begin{prop}\ollabel{cardplustimescommute}
$\cardplus$ et $\cardtimes$ sont commutatives et associatives.
\end{prop}

\begin{proof}
Pour la commutativité, d'après le
\olref[cardinals][cardsasords]{lem:CardinalsBehaveRight},
il suffit d'observer que
$\cardeq{(\cardfont{a} \disjointsum
\cardfont{b})}{(\cardfont{b} \disjointsum \cardfont{a})}$
et
$\cardeq{(\cardfont{a} \times \cardfont{b})}
{(\cardfont{b} \times \cardfont{a})}$.
Nous laissons l'associativité en exercice.
\end{proof}

\begin{prob}
Démontrez que $\cardplus$ et $\cardtimes$ sont associatives.
\end{prob}

\begin{prop}
$A$ est infini si et seulement si
$\card{A} \cardplus 1 = 1 \cardplus \card{A} = \card{A}$.
\end{prop}

\begin{proof}
Comme dans le
\olref[cardinals][classing]{generalinfinitycharacter},
cela découle du
\olref[ord-arithmetic][using-addition]{ordinfinitycharacter}
et du \olref[cardinals][cardsasords]{lem:CardinalsBehaveRight}.
\end{proof}

Nous voyons pourquoi il nous faut des symboles distincts
pour l'addition et la multiplication ordinales et cardinales :
il s'agit réellement d'opérations \emph{différentes}.
Les deux résultats suivants montrent que les deux
exponentiations diffèrent elles aussi.
(Rappelons que la \olref[z][infinity-again]{defnomega}
implique $2 = \{0, 1\}$.) 

\begin{lem}\ollabel{lem:SizePowerset2Exp}
$\card{\Pow{A}} = \cardexpo{2}{\card{A}}$,
pour tout ensemble $A$.
\end{lem}

\begin{proof}
Pour chaque sous-ensemble $B \subseteq A$,
définissons $\chi_B \in \funfromto{A}{2}$ par :
\begin{align*}
	\chi_{B}(x) &\defis
	\begin{cases}
		1 & \text{si }x\in B\\
		0 & \text{sinon.}
	\end{cases}
\end{align*}
Posons maintenant $f(B) = \chi_B$ ; cela définit
!!a{bijection} $f \colon \Pow{A}
\to \funfromto{A}{2}$. Donc
$\cardeq{\Pow{A}}{\funfromto{A}{2}}$.
Par suite,
$\cardeq{\Pow{A}}{\funfromto{\card{A}}{2}}$,
et ainsi
$\card{\Pow{A}} = \card{\funfromto{\card{A}}{2}}
= 2^{\card{A}}$.
\end{proof}

Cette démonstration concise englobe pour l'essentiel
la discussion de la \olref[sfr][siz][red-alt]{sec}.
Nous y avons montré comment «~ramener~» la
non-dénombrabilité de $\Pow{\omega}$ à celle de
l'ensemble des suites binaires infinies
$\Bin^\omega$. En effet, $\Bin^{\omega}$
n'est autre que $\funfromto{\omega}{2}$ ;
la démonstration précédente montre que le
raisonnement de la \olref[sfr][siz][red-alt]{sec}
reste valable si l'on remplace~$\omega$
par n'importe quel ensemble~$A$.
Ce résultat entraîne également un fait immédiat
sur l'exponentiation cardinale :

\begin{cor}\ollabel{cantorcor}
$\cardfont{a} < \cardexpo{2}{\cardfont{a}}$
pour tout cardinal~$\cardfont{a}$.
\end{cor}

\begin{proof}
Cela découle du théorème de Cantor
(\olref[sfr][siz][car]{thm:cantor}) et du
\olref{lem:SizePowerset2Exp}.
\end{proof}
\noindent
Ainsi, $\omega < \cardexpo{2}{\omega}$.
Attention toutefois : il s'agit ici de
l'exponentiation \emph{cardinale}. Pour
l'exponentiation \emph{ordinale}, au contraire,
$\omega = \ordexpo{2}{\omega}$
(voir la \olref[ord-arithmetic][expo]{sec}).

Puisque nous parlons d'exponentiation cardinale,
nous pouvons préciser de quelle «~manière~»
$\Real$ est !!{nonenumerable}.

\begin{thm}\ollabel{continuumis2aleph0}
$\card{\Real} = \cardexpo{2}{\omega}$
\end{thm}

\begin{proof}[Esquisse de démonstration]
Il existe plusieurs façons de le prouver. La plus
directe consiste à montrer que
$\cardle{\Pow{\omega}}{\Real}$ et
$\cardle{\Real}{\Pow{\omega}}$, puis à appliquer
le théorème de Schröder--Bernstein pour obtenir
$\cardeq{\Real}{\Pow{\omega}}$, et enfin
le \olref[card-arithmetic][opps]{lem:SizePowerset2Exp}
pour conclure que
$\card{\Real} = \cardexpo{2}{\omega}$.
Nous laissons comme exercice (instructif) la
construction d'injections
$f \colon \Pow{\omega} \to \Real$
et $g \colon \Real \to \Pow{\omega}$.
\end{proof}

\begin{prob}
Complétez la démonstration du
\olref[sth][card-arithmetic][opps]{continuumis2aleph0}
en montrant que
$\cardle{\Pow{\omega}}{\Real}$ et
$\cardle{\Real}{\Pow{\omega}}$.
\end{prob}

\end{document}
